\chapter{Étude et représentation des fonctions}
\textit{Tu sais désormais déterminer un domaine, calculer une limite et dériver une
fonction. Il est temps de réunir tous ces outils dans une seule démarche : l'\textbf{étude
complète} d'une fonction. Le but est de comprendre entièrement le comportement de $f$ —
ses variations, ses bornes, ses asymptotes — pour en tracer une courbe fidèle. Ce
chapitre est le point d'orgue de l'analyse de Première : maîtrise bien le \emph{plan
d'étude}, il revient à chaque devoir et à chaque épreuve.}

\section{Le plan d'étude d'une fonction}

Étudier une fonction $f$, c'est suivre une suite d'étapes ordonnées qui aboutissent à
un \textbf{tableau de variations} puis au \textbf{tracé} de sa courbe $\mathcal{C}_f$.

\begin{aretenir}[Les sept étapes de l'étude complète]
\begin{enumerate}
  \item \textbf{Ensemble de définition} $D_f$ : on exclut les valeurs interdites
        (dénominateur nul, radicande négatif).
  \item \textbf{Parité / symétries} : si $D_f$ est symétrique par rapport à $0$,
        tester $f(-x)$ ; on peut alors réduire l'intervalle d'étude.
  \item \textbf{Limites aux bornes} de $D_f$ (et de part et d'autre des valeurs exclues).
  \item \textbf{Dérivée} $f'$ et \textbf{signe} de $f'$.
  \item \textbf{Tableau de variations} : on y reporte signe de $f'$, sens de variation,
        valeurs et limites aux bornes.
  \item \textbf{Asymptotes et branches infinies} ; quelques \textbf{points particuliers}
        (intersections avec les axes, points remarquables).
  \item \textbf{Tracé} de $\mathcal{C}_f$ dans un repère.
\end{enumerate}
\end{aretenir}

\begin{remarque}
L'ordre n'est pas décoratif : chaque étape utilise la précédente. On ne dresse pas le
tableau de variations avant de connaître le signe de $f'$, et on ne trace pas la courbe
avant d'avoir le tableau et les asymptotes.
\end{remarque}

\section{Domaine, symétries et intervalle d'étude}

\begin{aretenir}[Réduire l'intervalle d'étude]
Si $f$ est \textbf{paire} ou \textbf{impaire} sur $D_f$ symétrique, il suffit d'étudier
$f$ sur $D_f\cap[0\,;\,+\infty[$ : on complète la courbe par symétrie (axe $(Oy)$ si $f$
est paire, centre $O$ si $f$ est impaire). On gagne ainsi la moitié du travail.
\end{aretenir}

\begin{exemple}
Pour $f(x)=\dfrac{x^2+1}{x^2-4}$, le domaine $D_f=\R\setminus\{-2\,;\,2\}$ est symétrique
par rapport à $0$ et $f(-x)=f(x)$ : $f$ est \textbf{paire}. On l'étudie sur
$[0\,;\,2[\,\cup\,]2\,;\,+\infty[$ et on complète par symétrie d'axe $(Oy)$.
\end{exemple}

\section{Limites aux bornes et continuité}

\begin{aretenir}
Aux bornes de $D_f$, on calcule les limites :
\begin{itemize}
  \item en $+\infty$ et $-\infty$ lorsque $D_f$ n'est pas borné ;
  \item à \textbf{gauche} et à \textbf{droite} de chaque valeur réelle exclue (valeur
        qui annule un dénominateur) ;
  \item au bord d'un intervalle fermé, on calcule simplement l'image.
\end{itemize}
Ces limites donnent le comportement de $\mathcal{C}_f$ aux extrémités et révèlent les
\textbf{asymptotes}.
\end{aretenir}

\begin{remarque}
Pour une fonction rationnelle $\dfrac{P(x)}{Q(x)}$, la limite en $\pm\infty$ s'obtient en
ne gardant que les \textbf{termes de plus haut degré} : $\displaystyle\lim_{x\to\pm\infty}
\frac{a_nx^n+\cdots}{b_px^p+\cdots}=\lim_{x\to\pm\infty}\frac{a_nx^n}{b_px^p}$.
\end{remarque}

\section{Dérivée, signe et tableau de variations}

\begin{theoreme}[Variations et signe de la dérivée]
Soit $f$ dérivable sur un intervalle $I$.
\begin{itemize}
  \item Si $f'(x)>0$ pour tout $x$ de $I$ (sauf en des points isolés), alors $f$ est
        \textbf{strictement croissante} sur $I$.
  \item Si $f'(x)<0$ pour tout $x$ de $I$, alors $f$ est \textbf{strictement décroissante}.
  \item Si $f'(x)=0$ pour tout $x$ de $I$, alors $f$ est \textbf{constante}.
\end{itemize}
\end{theoreme}

\begin{propriete}[Extremum local]
Si $f'$ s'\textbf{annule en changeant de signe} en $x_0$ (intérieur à $I$), alors $f$
présente un \textbf{extremum local} en $x_0$ :
\begin{itemize}
  \item un \textbf{maximum} si $f'$ passe du signe $+$ au signe $-$ ;
  \item un \textbf{minimum} si $f'$ passe du signe $-$ au signe $+$.
\end{itemize}
\end{propriete}

\begin{aretenir}[Dresser le tableau de variations]
Le tableau comporte trois lignes : la ligne des $x$ (avec $D_f$ et les valeurs où $f'$
s'annule), la ligne du \textbf{signe de $f'$}, et la ligne des \textbf{variations de $f$}
(flèches montantes/descendantes, valeurs et limites aux bornes). On y place une double
barre verticale sous chaque valeur exclue de $D_f$.
\end{aretenir}

\begin{exemple}
Pour $f(x)=x^3-3x$ sur $\R$, on a $f'(x)=3x^2-3=3(x-1)(x+1)$. Le signe de $f'$ est
positif sur $]-\infty\,;\,-1]$, négatif sur $[-1\,;\,1]$, positif sur $[1\,;\,+\infty[$.
\[
\begin{array}{c|ccccccc}
x        & -\infty &        & -1 &        & 1 &        & +\infty \\ \hline
f'(x)    &         & +      & 0  & -      & 0 & +      &         \\ \hline
         &         & \nearrow & 2 &        &   & \nearrow & +\infty \\
f        &-\infty  &        &    & \searrow & &        &         \\
         &         &        &    &        & -2 &       &
\end{array}
\]
Maximum local $f(-1)=2$, minimum local $f(1)=-2$.
\end{exemple}

\section{Asymptotes}

Les limites aux bornes décrivent comment la courbe « part à l'infini ». Lorsqu'elle
s'approche indéfiniment d'une droite, cette droite est une \textbf{asymptote}.

\begin{definition}[Asymptote verticale]
Si $\displaystyle\lim_{x\to a^-}f(x)=\pm\infty$ ou $\displaystyle\lim_{x\to a^+}f(x)=\pm\infty$,
la droite d'équation $x=a$ est une \textbf{asymptote verticale} à $\mathcal{C}_f$.
\end{definition}

\begin{definition}[Asymptote horizontale]
Si $\displaystyle\lim_{x\to+\infty}f(x)=\ell$ (réel), la droite d'équation $y=\ell$ est
une \textbf{asymptote horizontale} à $\mathcal{C}_f$ \textbf{en $+\infty$}. Définition
analogue en $-\infty$.
\end{definition}

\begin{definition}[Asymptote oblique]
La droite d'équation $y=ax+b$ (avec $a\neq 0$) est une \textbf{asymptote oblique} à
$\mathcal{C}_f$ en $+\infty$ si
\[ \lim_{x\to+\infty}\big[f(x)-(ax+b)\big]=0. \]
Définition analogue en $-\infty$.
\end{definition}

\begin{aretenir}[Rechercher une asymptote oblique]
\textbf{Méthode~1 (rationnelle simple).} Si $f(x)=\dfrac{P(x)}{Q(x)}$ avec
$\deg P=\deg Q+1$, effectuer la \textbf{division euclidienne} de $P$ par $Q$ :
\[ f(x)=ax+b+\frac{R(x)}{Q(x)},\qquad \text{avec } \frac{R(x)}{Q(x)}\xrightarrow[x\to\pm\infty]{}0. \]
La droite $y=ax+b$ est alors asymptote oblique.

\smallskip
\textbf{Méthode~2 (générale).} Calculer successivement
\[ a=\lim_{x\to+\infty}\frac{f(x)}{x},\qquad b=\lim_{x\to+\infty}\big[f(x)-ax\big]. \]
Si $a$ et $b$ sont finis (et $a\neq 0$), $y=ax+b$ est asymptote oblique.
\end{aretenir}

\begin{propriete}[Position de la courbe par rapport à une asymptote oblique]
On étudie le \textbf{signe de $f(x)-(ax+b)$}, qui vaut ici $\dfrac{R(x)}{Q(x)}$ :
\begin{itemize}
  \item si $f(x)-(ax+b)>0$, la courbe est \textbf{au-dessus} de l'asymptote ;
  \item si $f(x)-(ax+b)<0$, la courbe est \textbf{en dessous} ;
  \item le signe peut changer : la courbe peut \textbf{traverser} son asymptote.
\end{itemize}
\end{propriete}

\begin{exemple}
Soit $f(x)=\dfrac{x^2+1}{x}=x+\dfrac{1}{x}$. La division est déjà faite : comme
$\dfrac{1}{x}\to 0$ en $\pm\infty$, la droite $y=x$ est asymptote oblique. De plus
$f(x)-x=\dfrac{1}{x}$ : la courbe est \textbf{au-dessus} de l'asymptote pour $x>0$,
\textbf{en dessous} pour $x<0$. La droite $x=0$ est asymptote verticale.
\end{exemple}

\begin{illus}
\begin{tikzpicture}
\begin{axis}[
  width=10cm,height=8cm,
  axis lines=middle,
  xlabel={$x$},ylabel={$y$},
  xlabel style={right},ylabel style={above},
  xmin=-4.2,xmax=4.2,ymin=-5.2,ymax=5.2,
  xtick={-3,-2,-1,1,2,3},ytick={-4,-2,2,4},
  tick style={onbmuted},
  every axis x label/.style={at={(ticklabel* cs:1)},anchor=west},
  every axis y label/.style={at={(ticklabel* cs:1)},anchor=south},
  clip=true]
  % asymptote oblique y = x
  \addplot[onbmuted,dashed,thick,domain=-4:4]{x};
  \node[onbmuted] at (axis cs:3.2,4.4){$y=x$};
  % asymptote verticale x = 0 : l'axe Oy joue ce rôle (dashed pour rappel)
  % deux branches de f(x)=x+1/x
  \addplot[onbo,very thick,smooth,samples=100,domain=0.21:4]{x+1/x};
  \addplot[onbo,very thick,smooth,samples=100,domain=-4:-0.21]{x+1/x};
  \node[onbo] at (axis cs:2.6,4.6){$\mathcal{C}_f$};
\end{axis}
\end{tikzpicture}
\fig{Figure — courbe de $f(x)=x+\dfrac1x$ : asymptote verticale $x=0$ et asymptote oblique $y=x$ (en pointillés).}
\end{illus}

\section{Branches infinies}

Quand $\displaystyle\lim_{x\to+\infty}f(x)=\pm\infty$, la courbe possède une
\textbf{branche infinie} dont on précise la nature à l'aide de $\dfrac{f(x)}{x}$.

\begin{aretenir}[Nature d'une branche infinie en $+\infty$ (approche simple)]
On calcule $\displaystyle\lim_{x\to+\infty}\frac{f(x)}{x}$.
\begin{itemize}
  \item Si la limite est $0$ : branche \textbf{parabolique} de direction $(Ox)$.
  \item Si la limite est $+\infty$ (ou $-\infty$) : branche \textbf{parabolique} de
        direction $(Oy)$.
  \item Si la limite est un réel $a\neq 0$, on calcule $b=\lim[f(x)-ax]$ :
        \begin{itemize}
          \item $b$ fini : \textbf{asymptote oblique} $y=ax+b$ ;
          \item $b$ infini : branche parabolique de direction $y=ax$.
        \end{itemize}
\end{itemize}
\end{aretenir}

\begin{exemple}
Pour $f(x)=\sqrt{x^2+x}$ (sur $[0\,;\,+\infty[$), $\dfrac{f(x)}{x}=\sqrt{1+\frac1x}\to 1$,
puis $f(x)-x=\dfrac{x}{\sqrt{x^2+x}+x}\to\dfrac12$. La droite $y=x+\dfrac12$ est asymptote
oblique en $+\infty$.
\end{exemple}

\section{Points particuliers et tracé}

\begin{aretenir}[Avant de tracer]
On rassemble :
\begin{itemize}
  \item les \textbf{intersections avec les axes} : avec $(Oy)$ on calcule $f(0)$ (si
        $0\in D_f$) ; avec $(Ox)$ on résout $f(x)=0$ ;
  \item les \textbf{extremums} (lus dans le tableau de variations) ;
  \item les \textbf{asymptotes} et la \textbf{position} de la courbe par rapport à elles ;
  \item éventuellement le \textbf{coefficient directeur des tangentes} en quelques points
        (via $f'$), notamment aux points où la courbe coupe un axe.
\end{itemize}
On place ces éléments, puis on trace en respectant le sens de variation du tableau.
\end{aretenir}

\begin{remarque}
Une courbe bien tracée \emph{respecte le tableau} : elle monte là où $f'>0$, descend là
où $f'<0$, présente une tangente horizontale aux extremums, et longe les asymptotes
sans jamais s'en écarter à l'infini.
\end{remarque}

\section{Application aux fonctions usuelles}

\begin{propriete}[Fonctions polynômes]
Une fonction polynôme est définie sur $\R$, n'a \textbf{aucune asymptote} et ses
branches infinies sont \textbf{paraboliques} (de direction $(Oy)$). Sa dérivée est un
polynôme : son signe se lit en factorisant.
\end{propriete}

\begin{propriete}[Fonctions rationnelles]
$f=\dfrac{P}{Q}$ est définie là où $Q\neq 0$. Chaque racine de $Q$ (non racine de $P$)
donne une \textbf{asymptote verticale}. Selon les degrés :
\begin{itemize}
  \item $\deg P<\deg Q$ : asymptote horizontale $y=0$ ;
  \item $\deg P=\deg Q$ : asymptote horizontale $y=\dfrac{a_n}{b_n}$ (rapport des
        coefficients dominants) ;
  \item $\deg P=\deg Q+1$ : asymptote \textbf{oblique} (division euclidienne).
\end{itemize}
\end{propriete}

\begin{propriete}[Fonctions irrationnelles simples]
$f(x)=\sqrt{u(x)}$ est définie où $u(x)\ge 0$ ; elle a même sens de variation que $u$ là
où $u>0$, et $f'(x)=\dfrac{u'(x)}{2\sqrt{u(x)}}$. Ses branches infinies s'étudient via
$\dfrac{f(x)}{x}$ (souvent une asymptote oblique).
\end{propriete}

\begin{exemple}
$f(x)=\dfrac{2x-1}{x-1}=2+\dfrac{1}{x-1}$ : asymptote verticale $x=1$, asymptote
horizontale $y=2$ (degrés égaux, rapport $2/1$). La courbe est une hyperbole de centre
$\Omega(1\,;\,2)$.
\end{exemple}

\section{Position relative de deux courbes}

\begin{definition}[Position relative]
Soit $f$ et $g$ définies sur un même intervalle $I$, de courbes $\mathcal{C}_f$ et
$\mathcal{C}_g$. La \textbf{position relative} des deux courbes s'obtient en étudiant le
\textbf{signe de la différence} $d(x)=f(x)-g(x)$ :
\begin{itemize}
  \item $d(x)>0$ : $\mathcal{C}_f$ est \textbf{au-dessus} de $\mathcal{C}_g$ ;
  \item $d(x)<0$ : $\mathcal{C}_f$ est \textbf{en dessous} de $\mathcal{C}_g$ ;
  \item $d(x)=0$ : les courbes se \textbf{coupent} (point d'intersection).
\end{itemize}
\end{definition}

\begin{aretenir}[Méthode]
\textbf{1.} Former $d(x)=f(x)-g(x)$ et le \textbf{factoriser}. \textbf{2.} Dresser le
\textbf{tableau de signes} de $d(x)$. \textbf{3.} Conclure intervalle par intervalle ; les
racines de $d$ donnent les \textbf{abscisses des points d'intersection}.
\end{aretenir}

\begin{exemple}
Position de $\mathcal{C}_f$ ($f(x)=x^2$) par rapport à $\mathcal{C}_g$ ($g(x)=x$).
On forme $d(x)=x^2-x=x(x-1)$. Donc $d(x)>0$ sur $]-\infty\,;\,0[\cup]1\,;\,+\infty[$
($\mathcal{C}_f$ au-dessus), $d(x)<0$ sur $]0\,;\,1[$ ($\mathcal{C}_f$ en dessous), et les
courbes se coupent en $O(0\,;\,0)$ et $A(1\,;\,1)$.
\end{exemple}

\begin{illus}
\begin{tikzpicture}
\begin{axis}[
  width=9.5cm,height=7cm,
  axis lines=middle,
  xlabel={$x$},ylabel={$y$},
  xlabel style={right},ylabel style={above},
  xmin=-1.4,xmax=2.2,ymin=-0.8,ymax=3.4,
  xtick={-1,1,2},ytick={1,2,3},
  tick style={onbmuted},
  every axis x label/.style={at={(ticklabel* cs:1)},anchor=west},
  every axis y label/.style={at={(ticklabel* cs:1)},anchor=south},
  clip=true]
  \addplot[onbo,very thick,smooth,samples=80,domain=-1.2:1.85]{x^2};
  \node[onbo] at (axis cs:1.7,3.05){$\mathcal{C}_f:\,x^2$};
  \addplot[onbgreen,very thick,domain=-1.2:2.05]{x};
  \node[onbgreen] at (axis cs:1.95,1.55){$\mathcal{C}_g:\,x$};
  \filldraw[onbink] (axis cs:0,0) circle (2pt);
  \filldraw[onbink] (axis cs:1,1) circle (2pt) node[above left,onbink]{$A$};
\end{axis}
\end{tikzpicture}
\fig{Figure — position relative : $\mathcal{C}_f$ est en dessous de $\mathcal{C}_g$ sur $]0\,;\,1[$, au-dessus ailleurs ; intersections en $O$ et $A(1\,;\,1)$.}
\end{illus}

% ════════════════════════════════════════════════════════════
\section{L'essentiel à retenir}

\begin{synthese}
\textbf{Le plan d'étude en 7 étapes.}
\begin{enumerate}
  \item \textbf{$D_f$} : exclure dénominateurs nuls et radicandes négatifs.
  \item \textbf{Parité} (domaine symétrique) : $f(-x)=f(x)$ (paire, axe $Oy$) ou
        $f(-x)=-f(x)$ (impaire, centre $O$) $\Rightarrow$ réduire l'intervalle.
  \item \textbf{Limites} aux bornes et de part et d'autre des valeurs exclues.
  \item \textbf{Dérivée} $f'$ et son \textbf{signe} (factoriser).
  \item \textbf{Tableau de variations} : signe de $f'$ + flèches + valeurs/limites.
  \item \textbf{Asymptotes} et \textbf{points particuliers}.
  \item \textbf{Tracé} fidèle au tableau.
\end{enumerate}

\smallskip
\textbf{Asymptotes.}
\[ \lim_{x\to a}f=\pm\infty \Rightarrow x=a \text{ (verticale)};\quad
   \lim_{x\to\pm\infty}f=\ell \Rightarrow y=\ell \text{ (horizontale)}. \]
\textbf{Oblique} $y=ax+b$ : $\lim[f(x)-(ax+b)]=0$, ou division euclidienne si
$\deg P=\deg Q+1$, ou $a=\lim\frac{f(x)}{x}$ puis $b=\lim[f(x)-ax]$.

\smallskip
\textbf{Position / asymptote} : signe de $f(x)-(ax+b)$ (au-dessus si $>0$).

\smallskip
\textbf{Position de deux courbes} : signe de $d(x)=f(x)-g(x)$ ; racines = intersections.
\end{synthese}

% ════════════════════════════════════════════════════════════
\section{Méthodes \& savoir-faire}

\methode{Mener l'étude complète d'une fonction}
Dérouler les sept étapes dans l'ordre : domaine, parité, limites, dérivée et signe,
tableau de variations, asymptotes et points particuliers, tracé. Ne jamais sauter le
tableau de variations : c'est le pivot de toute l'étude.

\methode{Déterminer une asymptote horizontale ou verticale}
\textbf{Verticale} : repérer les valeurs $a$ exclues de $D_f$ et vérifier que
$\lim_{x\to a}f=\pm\infty$. \textbf{Horizontale} : calculer $\lim_{x\to\pm\infty}f$ ; si
la limite est un réel $\ell$, alors $y=\ell$ est asymptote.
\begin{exemple}
$f(x)=\dfrac{3x-2}{x+1}$. Valeur exclue $x=-1$ : $\lim_{x\to-1}f=\pm\infty$ donc $x=-1$
asymptote verticale. En $\pm\infty$, $f(x)\to\dfrac{3x}{x}=3$ : $y=3$ asymptote horizontale.
\end{exemple}

\methode{Trouver une asymptote oblique par division euclidienne}
Quand $\deg P=\deg Q+1$, poser la division de $P$ par $Q$ pour écrire
$f(x)=ax+b+\dfrac{R(x)}{Q(x)}$. La droite $y=ax+b$ est asymptote oblique, et le signe de
$\dfrac{R(x)}{Q(x)}$ donne la position de la courbe.
\begin{exemple}
$f(x)=\dfrac{x^2-x+1}{x-1}$. La division donne $f(x)=x+\dfrac{1}{x-1}$ (car
$x^2-x+1=x(x-1)+1$). Donc $y=x$ est asymptote oblique ; $f(x)-x=\dfrac{1}{x-1}>0$ pour
$x>1$ (courbe au-dessus), $<0$ pour $x<1$ (courbe en dessous).
\end{exemple}

\methode{Préciser une branche infinie par $\dfrac{f(x)}{x}$}
Lorsque $f(x)\to\pm\infty$, calculer $\lim\dfrac{f(x)}{x}$ : limite nulle (branche
parabolique direction $Ox$), infinie (direction $Oy$), ou réelle $a\neq 0$ (calculer
alors $b=\lim[f(x)-ax]$ pour conclure asymptote oblique ou branche parabolique).
\begin{exemple}
$f(x)=x+\sqrt{x}$ sur $[0\,;\,+\infty[$ : $\dfrac{f(x)}{x}=1+\dfrac{1}{\sqrt{x}}\to 1$,
puis $f(x)-x=\sqrt{x}\to+\infty$ : branche parabolique de direction la droite $y=x$.
\end{exemple}

\methode{Étudier la position de deux courbes}
Former $d(x)=f(x)-g(x)$, factoriser, dresser le tableau de signes ; conclure « au-dessus /
en dessous » et donner les abscisses d'intersection (racines de $d$).
\begin{exemple}
$f(x)=x^3$ et $g(x)=x$ : $d(x)=x^3-x=x(x-1)(x+1)$. La courbe $\mathcal{C}_f$ est en
dessous de $\mathcal{C}_g$ sur $]-\infty\,;\,-1[\cup]0\,;\,1[$ et au-dessus sur
$]-1\,;\,0[\cup]1\,;\,+\infty[$ ; intersections en $-1$, $0$ et $1$.
\end{exemple}

\methode{Exploiter parité et symétrie pour gagner du temps}
Tester la parité avant tout long calcul : si $f$ est paire ou impaire, n'étudier que
$[0\,;\,+\infty[$ et compléter le tracé par la symétrie correspondante.
\begin{exemple}
$f(x)=\dfrac{x}{x^2+1}$ est impaire (vu au chapitre 1) : on l'étudie sur $[0\,;\,+\infty[$
puis on complète $\mathcal{C}_f$ par symétrie de centre $O$.
\end{exemple}

\methode{Étude complète d'une fonction rationnelle (modèle)}
Étudions entièrement $f(x)=\dfrac{x^2+2}{x-1}$ en suivant les sept étapes.
\begin{exemple}
\textbf{1. Domaine.} $x-1\neq0$ : $D_f=\R\setminus\{1\}=\,]-\infty\,;\,1[\cup]1\,;\,+\infty[$.

\textbf{2. Parité.} $D_f$ n'est pas symétrique par rapport à $0$ : $f$ n'est ni paire ni
impaire. On étudie sur $D_f$ tout entier.

\textbf{3. Forme et limites.} Division euclidienne : $x^2+2=(x-1)(x+1)+3$, donc
\[ f(x)=x+1+\frac{3}{x-1}. \]
En $\pm\infty$ : $f(x)\sim x\to\pm\infty$. En $1$ : $\lim_{x\to1^-}\dfrac{3}{x-1}=-\infty$
donc $\lim_{x\to1^-}f=-\infty$ ; de même $\lim_{x\to1^+}f=+\infty$.

\textbf{4. Dérivée et signe.} $f'(x)=1-\dfrac{3}{(x-1)^2}=\dfrac{(x-1)^2-3}{(x-1)^2}$. Le
signe est celui de $(x-1)^2-3$, qui s'annule pour $x=1-\sqrt3$ et $x=1+\sqrt3$. Donc
$f'>0$ sur $]-\infty\,;\,1-\sqrt3[$ et $]1+\sqrt3\,;\,+\infty[$, et $f'<0$ entre les deux
(privé de $x=1$).

\textbf{5. Tableau de variations.} Maximum local en $x=1-\sqrt3$, minimum local en
$x=1+\sqrt3$ ; double barre en $x=1$.
\[
\begin{array}{c|ccccccccc}
x     & -\infty & & 1-\sqrt3 & & 1 & & 1+\sqrt3 & & +\infty \\ \hline
f'(x) & & + & 0 & - & \| & - & 0 & + & \\ \hline
f     & -\infty & \nearrow & & \searrow & \| & & \searrow & \nearrow & +\infty
\end{array}
\]

\textbf{6. Asymptotes et points particuliers.} $x=1$ est asymptote verticale ; $y=x+1$ est
asymptote oblique (car $\frac{3}{x-1}\to0$). Position : $f(x)-(x+1)=\dfrac{3}{x-1}$, donc
courbe \textbf{au-dessus} pour $x>1$, \textbf{en dessous} pour $x<1$. Intersection avec
$(Oy)$ : $f(0)=-2$.

\textbf{7. Tracé.} Deux branches d'hyperbole de centre $\Omega(1\,;\,2)$, longeant $x=1$ et
$y=x+1$, conformes au tableau de variations.
\end{exemple}
